\documentclass[12pt,a4paper]{article}
\usepackage[utf8]{inputenc}
\usepackage[russian]{babel}
\usepackage[T2A]{fontenc}
\usepackage{amsmath}
\usepackage{amssymb}
\usepackage{graphicx}
\usepackage{hyperref}
\usepackage{listings}
\usepackage{xcolor}
\usepackage{booktabs}
\usepackage{longtable}

\definecolor{codegreen}{rgb}{0,0.6,0}
\definecolor{codegray}{rgb}{0.5,0.5,0.5}
\definecolor{codepurple}{rgb}{0.58,0,0.82}
\definecolor{backcolour}{rgb}{0.95,0.95,0.92}

\lstdefinestyle{mystyle}{
    backgroundcolor=\color{backcolour},   
    commentstyle=\color{codegreen},
    keywordstyle=\color{magenta},
    numberstyle=\tiny\color{codegray},
    stringstyle=\color{codepurple},
    basicstyle=\ttfamily\footnotesize,
    breakatwhitespace=false,         
    breaklines=true,                 
    captionpos=b,                    
    keepspaces=true,                 
    numbers=left,                    
    numbersep=5pt,                  
    showspaces=false,                
    showstringspaces=false,
    showtabs=false,                  
    tabsize=2
}

\lstset{style=mystyle}

\title{Отчет №1: Реализация пайплайна оценки качества алгоритмов сравнения текстов}
\author{Команда разработки NIR 2026}
\date{\today}

\begin{document}

\maketitle

\tableofcontents
\newpage

\section{Введение}
В рамках разработки системы для текстологического анализа древнерусских текстов была поставлена задача автоматической оценки качества работы алгоритмов сравнения (diff algorithms). Для этого необходимо сопоставить результаты работы алгоритма с "эталонными" данными (ground truth), полученными из критического аппарата издания.

\section{Постановка задачи}
Цель: Разработать программный модуль и методику для автоматического сравнения вывода алгоритма с эталонными данными.

Входные данные:
\begin{enumerate}
    \item \textbf{Критический аппарат (.docx):} Файл, содержащий базовый текст (M10) и список разночтений для различных списков (сиглов, например P81, TB, MN). Разночтения представлены построчно.
    \item \textbf{Сырые тексты списков (.docx):} Файл, содержащий полные тексты фрагментов для каждой рукописи, но перемешанные по строкам.
\end{enumerate}

Выходные данные:
\begin{enumerate}
    \item DataFrame, сопоставляющий каждое слово базового текста с операцией из эталона и операцией из алгоритма.
    \item Метрики качества (в перспективе): точность (precision), полнота (recall) определения замен, вставок и удалений.
\end{enumerate}

\section{Реализация}

\subsection{Этап 1: Парсинг эталона (Parse Apparatus)}
Был разработан скрипт \texttt{parse\_apparatus.py} для извлечения эталонных операций (Word Ops) из файла критического аппарата.

\textbf{Логика работы:}
\begin{itemize}
    \item Скрипт читает файл \texttt{.docx}.
    \item Игнорирует техническую разметку (надстрочные знаки).
    \item Разбивает документ на блоки, соответствующие строкам базового текста M10.
    \item Парсит строки с вариантами, используя разделитель \texttt{]}.
    \item Для каждого слова базового текста определяет, есть ли для выбранного сигла (например, P81) изменение (replace, delete) или текст совпадает (equal).
\end{itemize}

\textbf{Результат:} Список \texttt{true\_values}, содержащий объекты операций с полями:
\begin{itemize}
    \item \texttt{op\_type}: тип операции (equal, replace, delete).
    \item \texttt{base\_token}: слово из базового текста.
    \item \texttt{other\_token}: слово из сравниваемого текста (или пусто при удалении).
\end{itemize}

\subsection{Этап 2: Подготовка текстов (Split Manuscripts)}
Для запуска алгоритма необходимо иметь чистые тексты рукописей. Был создан скрипт \texttt{split\_manuscripts.py}.

\textbf{Логика работы:}
\begin{itemize}
    \item Парсит файл, содержащий перемешанные фрагменты строк (формат: \texttt{Номер. Сигл: Текст}).
    \item Группирует фрагменты по сиглам.
    \item Сохраняет восстановленный текст каждой рукописи в отдельный файл (например, \texttt{data/clean/M10.docx}, \texttt{data/clean/P81.docx}).
\end{itemize}

\subsection{Этап 3: Запуск алгоритма (Algorithm Execution)}
В Jupyter Notebook (\texttt{testing\_quality.ipynb}) реализован запуск сравнения:
\begin{enumerate}
    \item Загрузка очищенных текстов M10 и P81.
    \item Токенизация (разбиение по пробелам).
    \item Вызов функции \texttt{algo.comparison.compare\_tokens} (используется алгоритм Майерса или Левенштейна).
    \item Получение списка \texttt{test\_values} (результат работы алгоритма).
\end{enumerate}

\subsection{Этап 4: Сопоставление результатов (Matching)}
Основная сложность заключалась в корректном сопоставлении списка эталонных операций и списка алгоритмических операций, так как тексты содержат повторяющиеся слова.

\textbf{Решение:}
Использована библиотека \texttt{pandas}.
\begin{enumerate}
    \item Списки преобразованы в DataFrame.
    \item Для каждого слова вычислен порядковый номер вхождения (\texttt{cumcount}) в рамках своего списка.
    \item Выполнен \texttt{Left Join} таблицы \texttt{true\_values} с таблицей \texttt{test\_values} по ключам \texttt{['base\_token', 'occurrence']}.
\end{enumerate}

Это позволяет однозначно сопоставить $i$-е вхождение слова "и" в эталоне с $i$-м вхождением слова "и" в результате алгоритма.

\begin{lstlisting}[language=Python, caption=Пример кода сопоставления]
df_true['occurrence'] = df_true.groupby('base_token').cumcount()
df_test['occurrence'] = df_test.groupby('base_token').cumcount()

merged_df = pd.merge(
    df_true,
    df_test,
    on=['base_token', 'occurrence'],
    how='left',
    suffixes=('_true', '_test')
)
\end{lstlisting}

\section{Результаты}
В результате работы пайплайна мы получаем таблицу, где для каждого слова базового текста можно сравнить:
\begin{itemize}
    \item \texttt{op\_type\_true}: Что должно было быть (по мнению эксперта).
    \item \texttt{op\_type\_test}: Что определил алгоритм.
\end{itemize}

Это позволяет автоматически выявлять ошибки алгоритма (например, ложные срабатывания замен или пропуски изменений) и численно оценивать качество его работы на реальных исторических текстах.

\end{document}
